\chapter{Datos y Metodología}

% --------------------------------------------------------
\section{Fuentes de datos}
% --------------------------------------------------------

\subsection{Variable objetivo: EMBI Ecuador}
% TODO: descripción de la serie, fuente, período 2004-2025, por qué se usa desde 2013

\subsection{Variables financieras externas}
% TODO: WTI, Gold, VIX, DXY, US Treasury 10Y, ETF EMB, ETF HYG — fuentes y justificación

\subsection{Variables macroeconómicas locales}
% TODO: reservas, balanza comercial, desempleo, liquidez M2 — frecuencia, fuente BCE/BM

\subsection{Datos NLP: GDELT v2}
% TODO: descripción de las 10 variables, filtro por mención a Ecuador, período desde sep-2013

% --------------------------------------------------------
\section{Preprocesamiento y construcción del dataset}
% --------------------------------------------------------

\subsection{Integración y alineación temporal}
% TODO: 3.047 días de trading, decisiones sobre NaN (exclusión de remesas y producción petrolera)

\subsection{División temporal del dataset}
% TODO: split 80/10/10 cronológico — justificación de no usar K-fold aleatorio
% Tabla: Train 2.406 días (2013-2024) | Val 274 días | Test 274 días

\subsection{Feature engineering}
% TODO: lags 1, 7, 30 días + rolling means 7 y 30 días → 132 features
% target_future = EMBI[t+1] como variable objetivo

\subsection{Selección de variables}
% TODO: Mutual Information sobre train exclusivamente → 60 variables seleccionadas
% Tabla de composición: 29 externas, 21 macro, 6 AR, 4 NLP

% --------------------------------------------------------
\section{Modelos y experimentos}
% --------------------------------------------------------

\subsection{Baseline: ARIMA}
% TODO: ARIMA(2,0,1) — selección por auto\_arima, validación rolling (ventana 3 años)

\subsection{XGBoost y Random Forest}
% TODO: RandomizedSearchCV, 100 iteraciones, TimeSeriesSplit 10 folds
% Hiperparámetros óptimos XGBoost (tabla)
% Wilcoxon Signed-Rank para comparar modelos

\subsection{Experimento de ablación}
% TODO: 4 grupos (A: AR, B: AR+Macro, C: AR+NLP, D: AR+Macro+NLP)
% Justificación: aislar contribución incremental de cada grupo de variables

\subsection{Causalidad de Granger}
% TODO: test sobre 4 variables NLP seleccionadas + avg\_tone + goldstein, lags 1-10

\subsection{Análisis de regímenes estructurales}
% TODO: Test de Chow (2020-03 y 2022-01), Quandt-Andrews, CUSUM, SHAP por subperiodo

\subsection{Autoencoder NLP}
% TODO: arquitectura 10→5→3→5→10, EarlyStopping, comparativa PCA-3 vs AE-3 vs original-4

% --------------------------------------------------------
\section{Métricas de evaluación}
% --------------------------------------------------------
% TODO: RMSE, MAE, R² — justificación de RMSE como métrica principal
% Nota sobre por qué el RMSE test es más relevante que el RMSE CV dado el alto std
